\begin{figure}[H]
\centering
\begin{tikzpicture}[x=0.94cm,y=0.85cm,>=Latex,font=\small]
 \begin{scope}
  \draw[->] (-3.5,0)--(3.5,0) node[right] {$u$};
  \draw[densely dashed] (-3.2,1.05)--(3.2,1.05);
  \draw[densely dashed] (-3.2,2.1)--(3.2,2.1);
  \node[left] at (-3.2,1.05) {$\beta/2$};
  \node[left] at (-3.2,2.1) {$\beta$};
  \draw[blue!65!black,very thick,domain=-3.2:3.2,samples=80]
      plot (\x,{1.98-0.36*exp(-2.3*\x*\x)});
  \node[blue!65!black,anchor=west] at (0.7,1.35) {$z(u)=u+i\beta\theta(u)$};
  \node[below] at (0,0) {$0$};
  \draw[<->] (0,1.62)--(0,2.1);
  \node[above,anchor=south west] at (0.05,2.1) {$d_0$};
  \node at (0,-0.65) {Positive mean on the curved $X$ contour};
 \end{scope}
 \begin{scope}[yshift=-4.2cm]
  \fill[orange!20] (-3,0.15) rectangle (-1.8,0.62);
  \fill[blue!12] (-1.8,0.15) rectangle (1.8,0.62);
  \fill[orange!20] (1.8,0.15) rectangle (3,0.62);
  \draw[->] (-3.5,0)--(3.5,0) node[right] {$u$};
  \foreach \x/\lab in {-3/-L_Y,-1.8/-L_X,0/0,1.8/L_X,3/L_Y} {
    \draw (\x,-0.07)--(\x,0.09);
    \node[below] at (\x,-0.08) {$\lab$};
  }
  \node at (0,0.385) {common interval};
  \node[above] at (-2.4,0.67) {shoulder};
  \node[above] at (2.4,0.67) {shoulder};
  \draw[thick,blue!65!black] (-1.8,1.3)--(1.8,1.3);
  \draw[thick,orange!75!black] (-3,1.8)--(3,1.8);
  \node[right] at (1.85,1.3) {$X$ wall};
  \node[right] at (3.05,1.8) {$Y$ wall};
 \end{scope}
\end{tikzpicture}
\caption{Two distinct geometries in the magnetic estimate. Above, the
main $X$ contour carries a positive mean; the central distance $d_0$
can become small when almost all $X$ sites are upper sites. Below,
the particle numbers set different logarithmic wall lengths. The sharp
trial field treats the common interval and its two shoulders separately.
The curves and lengths are schematic.}
\label{fig:magnetic-geometry}
\end{figure}
